\chapter{Допуск квартичного межрёберного инварианта и необходимость нового коннектора}

\section{Цель проверки}

Предыдущая глава закрыла прямую физическую вторую степень: при одном общем
Хиггсе смешанные произведения разных кромок равны нулю. Теперь необходимо
проверить, может ли тот же неизменённый носитель всё же породить
коэффициент-свободную межрёберную чувствительность через обычный
квартичный или более высокий односледовый полином.

Проверка обязана отличить три разных утверждения:
\begin{enumerate}
 \item на семейных матрицах можно записать скаляр, чувствительный к
 относительному кадру;
 \item этот скаляр правильно типизирован на физических кромках $u,d,e$;
 \item он возникает из одного уже принятого оператора и одного следа.
\end{enumerate}
Первое утверждение само по себе недостаточно.

\section{Полиномиальное следствие ортогональности кромок}

Пусть
\begin{equation}
 D=D_u+D_d+D_e,
 \qquad
 D_a=X_a\otimes T_a,
 \label{eq:v7-quartic-total-edge-operator}
\end{equation}
и для $a\ne b$
\begin{equation}
 D_aD_b^\dagger=D_a^\dagger D_b=0.
 \label{eq:v7-quartic-input-orthogonality}
\end{equation}
Тогда положительные операторы $D_a^\dagger D_a$ имеют попарно
ортогональные опоры; то же верно для $D_aD_a^\dagger$. Поэтому для любого
целого $n\ge1$
\begin{align}
 \Tr(D^\dagger D)^n
 &=\sum_a\Tr(D_a^\dagger D_a)^n,\label{eq:v7-source-moment-splitting}\\
 \Tr(DD^\dagger)^n
 &=\sum_a\Tr(D_aD_a^\dagger)^n.
 \label{eq:v7-target-moment-splitting}
\end{align}

Это не только квартичный no-go. Любой обычный односледовый спектральный
полином от текущего $D$ остаётся суммой трёх независимых кромочных
полиномов. Повышение степени не создаёт отсутствующий цикл.

\section{Квартичный след на коизометрическом вакууме}

На минимуме
\begin{equation}
 X_a=U_aV,
 \qquad
 U_a\in U(3),
 \qquad
 VV^\dagger=I_3,
 \qquad
 V^\dagger V=P_3.
 \label{eq:v7-quartic-coisometric-vacuum}
\end{equation}
Ранги физических поднятых кромок равны $9,9,3$. Поэтому
\begin{equation}
 \Tr(D^\dagger D)^2+\Tr(DD^\dagger)^2
 =2(9+9+3)=42
 \label{eq:v7-quartic-vacuum-trace}
\end{equation}
для любой точки $U(3)_u\times U(3)_d\times U(3)_e$.

Машинный аудит проверил моменты степеней $n=1,\ldots,6$ на общих
ненормированных полях. Максимальный относительный остаток разложения не
превышает
\begin{equation}
 3.93\cdot10^{-16}.
 \label{eq:v7-polynomial-splitting-residual}
\end{equation}
Квартичный след на случайных вакуумных кадрах менялся не более чем на
$8.53\cdot10^{-14}$ вокруг значения $42$.

\section{Family-only скаляры и ошибка типизации}

Если забыть физические метки кромок, можно построить относительную матрицу
\begin{equation}
 W_{ab}=X_aX_b^\dagger=U_aU_b^\dagger.
 \label{eq:v7-family-only-relative-matrix}
\end{equation}
Например, величины
\begin{equation}
 \Re\Tr W_{ab},
 \qquad
 \Re\Tr W_{ab}^2,
 \qquad
 |\Tr W_{ab}|^2
 \label{eq:v7-family-only-sensitive-scalars}
\end{equation}
действительно меняются вдоль относительных $U(3)$-орбит. В случайном
аудите диапазоны трёх величин равны соответственно
\begin{equation}
 3.733,
 \qquad
 4.492,
 \qquad
 4.141.
 \label{eq:v7-family-only-scalar-ranges}
\end{equation}

Но такое сокращение требует внедиагональной матричной единицы
$E_{ab}$ в пространстве меток $\Lambda_{\rm ch}$. Физическая классификация
дала
\begin{equation}
 \operatorname{End}_{\mathcal A_{\rm SM}-\mathcal A_{\rm SM}}
 (\Lambda_{\rm ch})\simeq\mathbb C^3,
 \label{eq:v7-quartic-edge-commutant}
\end{equation}
то есть разрешены только диагональные $E_{aa}$. Поэтому чувствительный
family-only скаляр не является физически типизированным инвариантом
текущего носителя.

\begin{nogocriterion}[Запрет забывания метки кромки]
Нельзя сначала удалить неэквивалентные бимодульные метки $u,d,e$, построить
из голых семейных матриц желаемый $W_{ab}$, а затем вернуть результат в
физическое действие. Внедиагональное сокращение меток является новым
коннектором, даже если его коэффициент положен равным единице.
\end{nogocriterion}

\section{Почему старый Krajewski rectangle не является готовым ответом}

Том IV уже предъявил настоящий положительный прецедент. Общий
$J$-парный четырёхрёберный rectangle дал
\begin{equation}
 \Tr D_{\rm rect}^4
 =104+16\cos\theta\,\Tr(P_-H_uH_d).
 \label{eq:v7-krajewski-cross-word-precedent}
\end{equation}
Здесь коэффициент смешанного слова действительно является кратностью
замкнутого пути, а не ручным порталом.

Однако тот оператор имеет четыре узла в первичном rectangle, восемь после
$J$-удвоения и размерность двадцать четыре на семейных блоках. Его
положительный результат явно предполагает ненулевой общий коннектор.
Текущий $H_{15}$-носитель содержит три попарно неэквивалентные кромки и не
содержит этот rectangle как подграф. Поэтому старую формулу нельзя перенести
без вывода нового узла или стрелки и повторной проверки gauge-, Real- и
order-one-условий.

\section{Граница спектрального кручения}

Спектральное кручение остаётся свидетелем того, что более богатый
функционал способен видеть $u$--$d$-данные. Но прежний строгий гейт уже
установил: его компоненты вычисляются из заданного $D_F$ и не выбирают
неизвестный $D_F$ без нового вариационного уровня. На текущем полностью
вырожденном коизометрическом вакууме обычный физический след тем более не
получает относительной чувствительности.

\begin{theorem}[Провал квартичного admission на текущем носителе]
На исправленном common-Higgs носителе $H_{15}$ любой обычный
односледовый полиномиальный момент текущего физического оператора
распадается в сумму независимых вкладов $u,d,e$. Квартичный момент равен
$42$ на всём вакууме $U(3)^3$ и не видит относительные кадры.
Чувствительные family-only скаляры существуют, но требуют запрещённого
внедиагонального сокращения неэквивалентных бимодульных меток. Поэтому
коэффициент-свободный квартичный межрёберный инвариант на неизменённом
носителе не допущен.
\end{theorem}

\section{Архитектурное следствие}

Следующий шаг уже нельзя формулировать как ``повысить степень''. Необходимо
вывести новый gauge- и Real-совместимый коннектор, который завершает
смешанный Krajewski-цикл и сам фиксирует межрёберную нормировку. Он должен:
\begin{enumerate}
 \item сохранять наблюдаемый пакет $H_{15}$ как физическое чтение;
 \item не добавлять произвольную матрицу смешивания;
 \item проходить order-zero, order-one, градуировку и Real-условие;
 \item порождать смешанное слово одним следом;
 \item не возвращать закрытый двойной семейный счёт.
\end{enumerate}

\begin{protocol}
\begin{enumerate}
 \item моменты степеней $1$--$6$: \textbf{точно распадаются по кромкам};
 \item максимальный относительный остаток: \textbf{$3.93\cdot10^{-16}$};
 \item прямой квартичный смешанный след: \textbf{ноль};
 \item квартичный след на вакууме: \textbf{$42$};
 \item зависимость от относительных $U(3)$-кадров: \textbf{нет};
 \item чувствительные family-only скаляры: \textbf{существуют};
 \item их физическая бимодульная типизация: \textbf{не проходит};
 \item старый Krajewski rectangle: \textbf{положительный прецедент};
 \item rectangle содержится в текущем $H_{15}$: \textbf{нет};
 \item повышение полиномиальной степени без коннектора:
 \textbf{закрыто отрицательно};
 \item следующий гейт: \textbf{минимальный совместимый с $H_{15}$
 смешанный коннектор}.
\end{enumerate}
\end{protocol}

Машинный сертификат сохранён в
\path{s2t_v7_quartic_cross_edge_invariant_admission_gate.py} и
\path{s2t_v7_quartic_cross_edge_invariant_admission_gate_results.json}.